\clearpage
\fancyhead[L]{\small BRIGHT SPOT ON THE MOON.}
\SourceLeaf{441}
\begin{center}
{\large [BRIGHT SPOT ON THE MOON.]}
\end{center}

\begin{center}
\noindent\rule{\linewidth}{0.4pt}

Göttingische gelehrte Anzeigen. No. 46. pp. 449..452. 22 March 1821.

\noindent\rule{\linewidth}{0.4pt}
\end{center}

We are all the more pleased to give here an extract from a letter that Dr. \textsc{Olbers} wrote to Councillor \textsc{Gauss} on 27 February about a phenomenon he observed on the dark part of the Moon’s surface on 5 February this year. The highly ingenious conjectures that this gifted astronomer adds about a phenomenon in itself as rare as it is remarkable deserve the closest attention and lead to an entirely new and fruitful way of looking at it.

On 5 February (writes Dr. \textsc{Olbers}) I saw the phenomenon on the Moon that has been called a lunar volcano. I wrote about it in my diary at once, in these very words. `On the fifth the sky was very clear, but there was already moonlight. On the dark part of the Moon I have never seen the phenomenon that has been taken for a burning volcano on the Moon as distinctly and strikingly as I did this evening. As usual, it appeared to be in Aristarchus. It was small, but quite \emph{strikingly} brighter than the rest of the Moon that was not lit by the Sun. It looked very much like a star and had exactly the appearance of a fixed star of the third magnitude northeast of the Moon.'

As it was cloudy on 6 February, Dr. \textsc{Olbers} did not look for the phenomenon any further. Soon afterward, however, English newspapers reported that Captain \textsc{Kater} had given the Royal Society in London an account on 7 February of a volcano he had seen on the Moon. Further observations had indeed convinced him that it was a volcano in eruption.

It seems, then (Dr. \textsc{Olbers} continues), that Mr. \textsc{Kater} saw the same phenomenon that caught my attention on 5 February, and went on observing it. I do not know what convinced him that it really was a burning volcano. But from all we know about the nature of the Moon and its very doubtful atmosphere, a burning volcano seems almost impossible. `I believe instead that the Earth actually appeared reflected in the flat, noticeably smooth, almost polished face of a large rock cliff belonging to Aristarchus.' The image of part of the Earth reflected in this way must have been far brighter than everything else lit only by the Earth. Everything else scatters earthshine in all directions, whereas the image is formed by light reflected in just \emph{one} direction. Even if that imperfect reflection sent back only one tenth of the earthshine (our actual mirrors reflect about half the light that falls on them) and the cliff face was only $2''$ in diameter, it could still appear as bright as a star of the sixth magnitude. --- On this view it is possible to explain, \emph{first}, why we always see the volcano-like phenomena only at certain places on the Moon. \emph{Second}, it is possible to explain why they come into view only rarely instead of in every lunation. The libration must be the same to within about $2^\circ$. --- There can hardly be any doubt that lunar cliffs might have faces that reflect light more or less like a mirror. On Earth there may be large glacier surfaces that can also reflect light like imperfect mirrors. I mention this only as an analogy, for glaciers on the Moon are just as
\SourceLeaf{442}%
unlikely as burning volcanoes. But it seems quite conceivable that some of the mountains, probably formed on the Moon too according to laws of crystallization, could have flat, smooth, almost polished faces. There may be many of these on the Moon, but they may rarely be positioned just right to reflect the image of the Earth back to us at a given libration. The Sun, too, sometimes seems to be reflected in similar cliff faces on the Moon. As recently as about 8 weeks ago I saw two sunlit mountain peaks in Mare Imbrium beyond the terminator. They were so unusually bright, scintillating and Sirius-like that I could not suppose the sunlight was merely being reflected according to the ordinary laws of scattering. I cannot say for certain whether one of these mountains was perhaps Lahire (after \textsc{Schröter}), also in Mare Imbrium, where \textsc{Schröter} observed very similar phenomena.

In a later letter Dr. \textsc{Olbers} adds that this hypothesis about the cause of the phenomenon will be easy to test. If it gives the true explanation, the same phenomenon must always recur at the same libration. Cliff faces that can reflect an image of the Earth or the Sun more or less imperfectly are all the more conceivable on the Moon, he says, because the outer surfaces of the mountains and cliffs there probably do not weather through atmospheric action as they do on Earth. The reflecting cliff face need not be entirely flat, provided only that the reflecting parts lie in parallel planes, as could easily happen in mountains formed according to laws of crystallization. Here Dr. \textsc{Olbers} points to our basalt mountains. Their individual large crystals could very well still act together to reflect an imperfect image of the Sun to a distant eye, if their surfaces had not long ago lost their probably original polish and smoothness through air, vapors, rain, and so on.

On 6 March, when the night side of the Moon could be seen exceptionally well, Dr. \textsc{Olbers} could distinguish all the spots very clearly with his \textsc{Dollond} telescope, for example Grimaldi, Copernicus, Kepler, Maeilius, Menelaus, and so on. Aristarchus again stood out from all the others, apparently even \emph{more than usual}. This time, however, Dr. \textsc{Olbers} found it less \emph{bright} and less \emph{like a fixed star} than on 5 February.
